\section{Stability above a strict barrier}
\label{sec:big-reference-stability}

The subsolutions of \cref{lem:big-reference-subsolution} stay within a bounded distance of a given reference. With them, the proof of \cref{thm:bounded-one-sided-stability} extends to unbounded potentials lying above $\rho$, and the resulting estimate is uniform over all references between $\rho$ and zero. For a fixed reference, it upgrades $L^1$ compactness of solutions to uniform convergence of their differences on $\Omega$.

\begin{theorem}\label{thm:big-reference-one-sided-stability}\index{Stability!varying references}
Fix $C_0,L>0$, $p>1$, and
\[
 0<\gamma<\frac{p-1}{np+p-1}.
\]
Let $\eta\in\PSH(X,\theta)$ satisfy $\rho\leq\eta\leq0$, and suppose that $u\in\PSH(X,\theta)$ satisfies
\[
 \eta-C_0\leq u\leq\eta,\qquad
 \mathrm{MA}_\theta(u)=f\,\mathrm{d}V,\qquad 0\leq f\in L^p(\mathrm{d}V),
 \quad\|f\|_p\leq L.
\]
There is a constant $C$, depending only on the fixed geometry and on $\rho,C_0,L,p,\gamma$, such that every $v\in\PSH(X,\theta)$ with $v\leq\eta$ satisfies
\begin{equation}\label{eq:big-reference-stability}
 \sup_\Omega(v-u)_+
 \leq C\left(\int_\Omega(v-u)_+\,\mathrm{d}V\right)^\gamma.
\end{equation}
\end{theorem}

The comparison potential $v$ need not have the ordinary singularity type of $\eta$. Since $Z$ is pluripolar and has zero volume, the integral may equally be taken over $X$, and the resulting inequality $v\leq u+C\bigl(\int_\Omega(v-u)_+\,\mathrm{d}V\bigr)^\gamma$ on $\Omega$ extends to $X$.

\begin{proof}
Replacing $v$ by $\max(v,u)$, we may assume $\eta-C_0\leq u\leq v\leq\eta$, so that $0\leq v-u\leq C_0$ on $\Omega$. We repeat the proof of \cref{thm:bounded-one-sided-stability}, with $1<q<p$, $E_\delta=\{u<v-\delta\}$, $g_\delta=\mathbf1_{E_\delta}f$ and $m_\delta=\|g_\delta\|_q$, and with suprema taken over $\Omega$. Three ingredients change. First, when $m_\delta>0$, \cref{lem:big-reference-subsolution}, applied to $g_\delta/m_\delta$ and followed by subtracting one from its potential, gives $h\in\PSH(X,\theta)$ and a fixed $b>0$ such that
\[
 \eta-1\leq h\leq\eta,\qquad
 \mathrm{MA}_\theta(h)\geq b\,g_\delta/m_\delta\,\mathrm{d}V;
\]
the constant $b$ is uniform in $\eta$, because the auxiliary density has $L^q$ norm one and the strict barrier is fixed. These bounds give $h-v\leq C_0$ and $v-h\leq1$, the two properties of $h$ used there. Second, all potentials involved are bounded below by $\rho$ minus a constant, so positivity of mixed products holds on $\Omega$, where they are locally bounded, and $Z$ is not charged. Third, bounded domination is replaced by \cref{thm:big-barrier-domination}: the potentials $v-\delta$ and $w=(1-t)v+th-\delta-C_0t$ lie below $u+C_0$, so this theorem applies with the coefficients zero and $1/2$. The rest of the argument is unchanged. It gives \eqref{eq:big-reference-stability} with the exponent $\gamma_q=(p-q)/(npq+p-q)$, and we choose $q>1$ so close to one that $\gamma_q>\gamma$.
\end{proof}

\Needspace{15\baselineskip}
\begin{corollary}\label{cor:big-reference-uniform-convergence}
Fix $\eta\in\PSH(X,\theta)$ with $\rho\leq\eta\leq0$, constants $C_0,L>0$, and $p>1$. Consider the family of potentials $u\in\PSH(X,\theta)$ satisfying
\[
 \eta-C_0\leq u\leq\eta,\qquad
 \mathrm{MA}_\theta(u)=f_u\,\mathrm{d}V,\qquad0\leq f_u\in L^p(\mathrm{d}V),
 \quad\|f_u\|_p\leq L.
\]
This family is compact for uniform convergence of differences on $\Omega$. If $u_j$ belong to the family and $u_j\to u$ in $L^1$, then $u$ belongs to the family and
\[
 \sup_\Omega|u_j-u|\leq C\|u_j-u\|_{L^1(\mathrm{d}V)}^\gamma
 \longrightarrow0
\]
for every exponent $\gamma$ as in \cref{thm:big-reference-one-sided-stability}. Moreover, $f_{u_j}$ converges weakly to $f_u$ in $L^p$.
\end{corollary}
\begin{proof}
Applying the theorem in both orders to $u_j$ and $u_k$, with the common reference $\eta$, we see that an $L^1$-Cauchy sequence in the family is uniformly Cauchy in differences on $\Omega$. Its uniform limit on $\Omega$ is locally $\theta$-plurisubharmonic and agrees almost everywhere, hence everywhere for the canonical representatives, with the $L^1$ limit $u$. The bounds $\eta-C_0\leq u\leq\eta$ hold on all of $X$. Letting $k\to\infty$ in the pairwise estimate proves the quantitative bound, before any equation for $u$ is used.

The densities form a bounded set in the reflexive space $L^p$. Let $f$ be the weak limit of a subsequence $f_{u_{j_k}}$. Since the potentials are locally bounded and converge locally uniformly on $\Omega$, Bedford--Taylor continuity, which holds also for the smooth nonclosed background, gives $\mathrm{MA}_\theta(u)=f\,\mathrm{d}V$ on $\Omega$. Neither measure charges the analytic set $Z$, so this equality holds on $X$. Positivity and weak lower semicontinuity of the norm give $f\geq0$ and $\|f\|_p\leq L$. Thus every weak cluster point equals the density $f_u$ of $\mathrm{MA}_\theta(u)$, which proves weak convergence of the whole sequence. Finally, the common bounds by $\eta$ and $\eta-C_0$ give $L^1$ compactness of the family, by the compactness of quasi-plurisubharmonic functions, and the preceding argument upgrades each convergent subsequence to uniform convergence of differences on $\Omega$.
\end{proof}

For solutions with minimal singularities, normalized by $\sup_Xu=0$, we take $\eta=V_\theta$. Then $\rho\leq V_\theta\leq0$, and every normalized comparison potential $v\in\PSH(X,\theta)$ satisfies $v\leq V_\theta$, as the theorem requires.
